\documentclass[12pt]{article}
\usepackage[T1]{fontenc}
\usepackage[utf8]{inputenc}
\usepackage{lmodern}
\usepackage{textcomp}
\usepackage{enumitem}
\usepackage{geometry}
\geometry{margin=1in}
\begin{document}
\begin{center}
Answer Key - Constitutional Law - Graduate Level Assessment 15 \
\small (Answers are ordered exactly as in the question paper.)
\end{center}

\section*{Multiple Choice}
\begin{enumerate}[label=\arabic*.]
\item \textbf{Answer:} B
\\ \textit{Options:} A) Law is unstable.; B) Law is determinate.; C) Law reflects economic power.; D) Law is politics.; E) Law is a social construct.
\\ \textit{Note:} Correct option: B - Law is determinate.
\item \textbf{Answer:} A
\\ \textit{Options:} A) He claims that by giving priority to the needs of the poor, we can increase equality.; B) He rejects the idea of equality altogether.; C) He believes that equality is not a realistic goal.; D) He argues than an unequal society is inevitable.; E) He suggests that societal structures inherently support inequality.
\\ \textit{Note:} Correct option: A - He claims that by giving priority to the needs of the poor, we can increase equality.
\item \textbf{Answer:} A
\\ \textit{Options:} A) Legal Pragmatism; B) Legal Formalism; C) Comparative; D) Analytical; E) Sociological
\\ \textit{Note:} Correct option: A - Legal Pragmatism
\item \textbf{Answer:} C
\\ \textit{Options:} A) Because the concept of punishment has evolved over time.; B) Because punishment can be justified in multiple ways.; C) Because any definition of punishment should be value-neutral.; D) Because the practice of punishment is separate from its justification.; E) Because the justification of punishment varies across cultures.
\\ \textit{Note:} Correct option: C - Because any definition of punishment should be value-neutral.
\item \textbf{Answer:} B
\\ \textit{Options:} A) Execution of the arrest warrant; B) investigative surveillance; C) Police interrogation prior to arrest; D) Post-charge lineups; E) Preliminary Hearing
\\ \textit{Note:} Correct option: B - investigative surveillance
\end{enumerate}

\section*{True / False}
\begin{enumerate}[label=\arabic*.]
\setcounter{enumi}{5}
\item \textbf{Answer:} True
\\ \textit{Options:} A) True; B) False
\\ \textit{Note:} LLM-generated true statement based on MCQ.
\item \textbf{Answer:} True
\\ \textit{Options:} A) True; B) False
\\ \textit{Note:} LLM-generated true statement based on MCQ.
\item \textbf{Answer:} True
\\ \textit{Options:} A) True; B) False
\\ \textit{Note:} LLM-generated true statement based on MCQ.
\item \textbf{Answer:} True
\\ \textit{Options:} A) True; B) False
\\ \textit{Note:} LLM-generated true statement based on MCQ.
\item \textbf{Answer:} True
\\ \textit{Options:} A) True; B) False
\\ \textit{Note:} LLM-generated true statement based on MCQ.
\end{enumerate}

\section*{Long Form Response}
\begin{enumerate}[label=\arabic*.]
\setcounter{enumi}{10}
\item \textbf{Answer:} A diary entry from the insured's mother noting the day of the insured's birth.. Explain why this is the correct answer and provide supporting evidence. Reference key facts or concepts that support this choice.
\\ \textit{Note:} Long-form derived from MCQ; award credit for stating the correct idea and giving supporting reasoning.
\item \textbf{Answer:} It regards morals and law as inseparable.. Explain why this is the correct answer and provide supporting evidence. Reference key facts or concepts that support this choice.
\\ \textit{Note:} Long-form derived from MCQ; award credit for stating the correct idea and giving supporting reasoning.
\end{enumerate}

\vfill
\noindent\textit{End of Answer Key}
\end{document}